% -------------------------------------------------------
%  Abstract
% -------------------------------------------------------

\شروع{وسط‌چین}
\مهم{چکیده}
\پایان{وسط‌چین}
\بدون‌تورفتگی

انطباق در زمان آزمون برخط تلاش می‌کند تا تعمیم‌پذیری شبکه‌های عصبی عمیق از پیش آموزش‌دیده در مواجهه با داده‌های تغییر توزیع یافته و ناشناخته در زمان اجرا را بهبود ببخشد. در این فرآیند، مدل‌ها به‌صورت پویا و متناسب با داده‌های آزمون تنظیم می‌شوند. چالش اصلی این رویکرد، عدم دسترسی به کل مجموعه‌داده‌ی آزمون، فقدان برچسب برای نمونه‌ها و همچنین عدم امکان تنظیم دقیق فراپارامترهای شبکه به‌منظور دستیابی به بهترین عملکرد است. در چنین شرایطی، در صورت نبود کنترل‌های مناسب، فرایند انطباق می‌تواند منجر به انباشت خطا و پدیده‌ی فراموشی فاجعه‌بار شود.

برای کاهش این مشکلات، روش‌های موجود معمولاً از معیارهایی همچون بیشینه‌ی خروجی لایه‌ی پیش‌بینی مانند اندازه‌ی بیشینه و یا آنتروپی به‌عنوان شاخص اعتماد استفاده کرده‌اند. هدف این معیارها شناسایی نمونه‌هایی است که احتمال کمتری برای بروز خطا دارند. با این حال، مشاهدات نشان می‌دهند که این شاخص‌ها هرچند در داده‌های هم‌توزیع در زمان آموزش عملکرد مناسبی دارند، اما در شرایطی که توزیع داده‌ها در زمان آزمون تغییر کند، به‌ویژه در ابعاد بالا، دقت آن‌ها به‌شدت کاهش می‌یابد.

بر این اساس، روشی نوین برای انطباق در زمان آزمون پیشنهاد می‌شود که مسئله را از منظر شناسایی نمونه‌های خارج از توزیع بررسی کرده و با هدف دستیابی به استحکام، قدرت انطباق مدل را نیز حفظ می‌کند. روش پیشنهادی ساده، کارآمد و بدون نیاز به تنظیم فراپارامتر است که آن را برای کاربردهای مختلف در یادگیری برخط به گزینه‌ای ایده‌آل تبدیل می‌سازد. 

به‌طور‌خاص در مجموعه‌داده‌ی \lr{CIFAR100-C} نتایج تجربی نشان می‌دهند که روش پیشنهادی به‌طور متوسط حدود ۱۵٪ عملکرد بهتری نسبت به روش‌های پایه ارائه کرده است. در حالتهای آزمون مستمر که مدل‌های پایه دچار فروپاشی می‌شوند، روش پیشنهادی دچار این مشکل نشده و دقت ۷۵٪ را ارئه می‌دهد که نزدیک نتایج در حالت انطباق در زمان آزمون عادی می‌باشد. علاوه بر این، تحلیل روی مجموعه‌ای از فراپارامترهای مختلف نشان می‌دهد که میانگین عملکرد روش پیشنهادی همواره به مقدار بهینه نزدیک‌تر است، در حالی که روش‌های پایه نسبت به انتخاب فراپارامتر بسیار حساس هستند. به عنوان نمونه، بیشینه‌ی دقت روش پایه‌ای \lr{Tent} در برخی تنظیمات ۷۰٪ بوده، در حالی که نسخه‌ی بهبود یافته به ۷۹٪ رسیده است. ازطرفی میانگین دقت روش پیشنهادی برابر با ۷۸٪ و نزدیک به مقدار بهینه باقی مانده، اما میانگین روش پایه به ۵۹٪ کاهش یافته است. این یافته‌ها گویای پایداری و استحکام بالاتر روش پیشنهادی در شرایط تغییر توزیع و تنظیمات مختلف هستند.

\پرش‌بلند
\بدون‌تورفتگی \مهم{کلیدواژه‌ها}: 
انطباق زمان آزمون، یادگیری مستمر، شناسایی خارج از توزیع
\صفحه‌جدید
